%------------------------------------------------------------------------------%
\begin{question}
  Considere os pontos
  \[
  A=(0,0,0)
  \]
  \[
  B=(2,1,0)
  \]
  \[
  C=(1,3,2)
  \]
  \[
  D=(0,1,4)
  \]
  Determine o volume do tetraedro $ABCD$
\end{question}

%------------------------------------------------------------------------------%
\begin{resolution}{Exemplo 4 — Resolução}
Construímos
\[
\overrightarrow{AB}=(2,1,0)
\]
\[
\overrightarrow{AC}=(1,3,2)
\]
\[
\overrightarrow{AD}=(0,1,4)
\]
O paralelepípedo correspondente possui volume
\[
V_P=
\left|
\begin{vmatrix}
2&1&0\\
1&3&2\\
0&1&4
\end{vmatrix}
\right|
\]
\end{resolution}

%------------------------------------------------------------------------------%
\begin{resolution}{Exemplo 4 — Resolução}
Calculando o determinante
\[
V_P=
\left|
2
\begin{vmatrix}
3&2\\
1&4
\end{vmatrix}
-
\begin{vmatrix}
1&2\\
0&4
\end{vmatrix}
\right|
\]
\[
V_P=|2(10)-4|
\]
\[
V_P=16
\]
Um tetraedro ocupa um sexto do paralelepípedo determinado pelos mesmos três vetores
\[
V_T=\frac16V_P
\]
\[
V_T=\frac{16}{6}
\]
\[
V_T=\frac83
\]
\end{resolution}
%------------------------------------------------------------------------------%
